% ============================================================================
% SECTION 2.5 : LA FAMILLE SOTP — CONSTRUCTIONS OPTIMALES
% Source : partie2/sotp_construction.tex
% Version : 2.2 — alignée sur la formalisation Lean4 (section~\ref{sec:lean4})
% ============================================================================

\section{LA FAMILLE SOTP~: UNE CONSTRUCTION PAR R\'EGIME}
\label{sec:sotp_construction}
\begin{sloppypar}
Le Trilemme CRO conditionnel classique~(section~\ref{sec:cro_classique}) \'etablit
que $\Conf \cdot \Rel \cdot \Opp \leq \Delta/(\gamma\lambda) +
O(\lambda^{-2})$, avec une forme raffin\'ee. Une borne d'impossibilit\'e sans
construction est incompl\`ete, car elle \'enonce une limite sans dire si
cette limite est approchable. La \textbf{famille SOTP} (Semantic One-Time
Pad) propose des constructions de r\'ef\'erence, une par r\'egime de richesse
contextuelle, dont les propri\'et\'es cryptographiques sont prouv\'ees et dont
l'optimalit\'e \`a facteur constant est formul\'ee comme conjecture
(Conjecture~\ref{conj:optimalite_universelle}). Cette mise en regard
borne/construction transforme un r\'esultat n\'egatif en guide de conception.

\medskip
\noindent\textbf{Statut de cette contribution.} La famille SOTP
constitue la contribution \textbf{C4} de cette th\`ese. Elle est
aussi importante que la borne elle-m\^eme, puisqu'elle
donne au Trilemme CRO sa port\'ee op\'erationnelle. Les trois variantes
(SOTP-Low pour le contexte rare, SOTP-Base pour la transition, SOTP-High
pour le contexte riche) couvrent l'ensemble du spectre
d'entropie contextuelle et fournissent chacune une construction
cryptographique concr\`ete, dont la saturation de la borne reste \`a d\'emontrer. Cette sym\'etrie
se retrouve dans les th\'eor\`emes d'impossibilit\'e f\'econds. CAP trouve
son sens dans les syst\`emes distribu\'es r\'eels et FLP dans les protocoles
partiels (Paxos, Raft). Le Trilemme CRO trouve le sien dans les trois
variantes SOTP.

Cette section correspond \`a l'\'etape \textbf{S (Saturer)} de la progression
F$\to$B$\to$E$\to$\textbf{S}$\to$C$\to$X et apporte la contribution \textbf{C4}
de cette th\`ese.

La philosophie architecturale
commune est d'abord pos\'ee en sous-section~\ref{subsec:philosophie_sotp}. Les trois
constructions sont d\'etaill\'ees ensuite par ordre croissant de complexit\'e
contextuelle, avec SOTP-Low pour contexte rare (sous-section~\ref{subsec:sotp_low}),
SOTP-Base pour le r\'egime de transition (sous-section~\ref{subsec:sotp_base}),
SOTP-High pour contexte riche (sous-section~\ref{subsec:sotp_high}). Les
crit\`eres de s\'election entre variantes pour guider la conception
d'architecture sont fournis en sous-section~\ref{subsec:synthese_selection}.
La section se conclut par la conjecture d'optimalit\'e universelle et
l'analyse de sensibilit\'e (sous-section~\ref{subsec:optimalite_sensibilite}).
\end{sloppypar}

\begin{boiteArbitre}
La borne existe sans \^etre pour autant inaccessible. L'arbitre dispose
d\'esormais d'un protocole con\c{c}u \emph{pour lui}, \`a savoir SOTP-Low quand son
contexte est rare, SOTP-Base quand il est interm\'ediaire, SOTP-High quand il
est riche. Chaque variante calibre ses couches Gold,
Iron, Clay sur ce que son $\Delta$ lui permet d'extraire,
sans lui en demander davantage. L'impossibilit\'e devient ainsi un guide de
conception. Savoir \`a quel point ces constructions approchent la borne reste
une question ouverte. La borne sert d\'esormais de carte au concepteur.
\end{boiteArbitre}

% -----------------------------------------------------------------------
\subsection{Philosophie de conception~: paradigme Gold-Iron-Clay}
\label{subsec:philosophie_sotp}
% -----------------------------------------------------------------------
\begin{sloppypar}
Nous d\'efinissons l'architecture commune \`a toutes les variantes SOTP.

Trois couches fonctionnelles structurent chaque protocole selon le paradigme
Gold-Iron-Clay~(sous-section~\ref{subsec:gold_iron_clay}).
La couche \textbf{Gold} accomplit le masquage cryptographique au moyen de one-time pads
pseudo-al\'eatoires dans les r\'egimes bas et interm\'ediaire~\cite{Shannon49,BlumMicali84}, preuves zero-knowledge dans
le r\'egime riche~\cite{GMR85,BenOr88}. La couche \textbf{Iron} assure le binding via des commitments
contraignants, Pedersen~\cite{Pedersen91,ChaumPedersen92} pour SOTP-Low/Base, Groth16~\cite{Groth16,BCTV14}
pour SOTP-High, permettant l'ouverture v\'erifiable du message. La couche
\textbf{Clay} enfin pr\'esente le contexte institutionnel sous une forme que des
interpr\`etes \`a capacit\'e born\'ee~$\Delta$ peuvent exploiter pour rendre leurs
verdicts~\cite{Simon82,Miller56}~; elle n'accro\^it $\Opp$ que si elle transmet une
information que l'institution ne d\'etient pas d\'ej\`a (Remarque~\ref{rem:statut_sotp}).

L'architecture elle-m\^eme \'etait d\'ej\`a esquiss\'ee par la
construction S-OTP de la section~\ref{subsec:construction_sotp}. L'apport propre \`a cette section tient \`a son adaptation
par r\'egime, puisque chaque variante SOTP sp\'ecialise la couche Clay selon l'entropie
contextuelle $H(\mathcal{C})$, substituant encodage d\'eterministe, compression
Johnson-Lindenstrauss~\cite{JohnsonLindenstrauss84,DasguptaGupta03,Vempala05}, ou hachage SNARK
Poseidon~\cite{Groth16,Poseidon2019} selon la richesse du contexte.

\subsubsection{Correspondance entre m\'etriques CLO et propri\'et\'es cryptographiques}
\label{subsubsec:interpretation_operationnelle}

Tandis que le cadre CLO d\'efinit $\Conf$, $\Rel$, $\Opp$ via la th\'eorie de
l'information~\cite{MINKA2026CLO,Shannon48,CoverThomas06}, les preuves de ce chapitre les
expriment en termes de propri\'et\'es de s\'ecurit\'e cryptographique standard admettant des
r\'eductions aux hypoth\`eses bien \'etablies~\cite{Goldreich01,KatzLindell14}.

La \textbf{confidentialit\'e} $\Conf$, notion statistique, a pour pendant calculatoire l'indistinguabilit\'e s\'emantique~\cite{GoldwasserMicali84,MicaliRackoff91}, selon laquelle, pour
toute paire de messages $(m_0, m_1)$ dans le m\^eme contexte $c$, nul adversaire polynomial
ne peut distinguer la preuve de $m_0$ de celle de $m_1$ avec avantage non n\'egligeable.
La \textbf{fiabilit\'e de v\'erification} $\Relver$ (Remarque~\ref{rem:rel_ver}) correspond au binding~\cite{Pedersen91,Damgard87}, propri\'et\'e selon laquelle aucun adversaire polynomial ne peut
produire deux ouvertures distinctes d'un m\^eme commitment sur des messages diff\'erents avec
probabilit\'e non n\'egligeable. Quant \`a l'\textbf{opposabilit\'e} $\Opp$, elle correspond \`a
la fraction d'information extractible sur le verdict, soit $I(V;\Sigma|C)/H(V|C)$, quantifi\'ee
via la capacit\'e d'extraction institutionnelle~$\Delta$ (D\'efinition~\ref{def:delta_extraction})~\cite{Simon82,Miller56,Ericsson95}.

Cette traduction rend le cycle des preuves explicite, car les garanties de confidentialit\'e
et de fiabilit\'e de v\'erification r\'esultent de r\'eductions formelles \`a des hypoth\`eses
cryptographiques standard (PRG~\cite{BlumMicali84,Yao82}, DLog~\cite{DiffieHellman76},
q-SDH/q-PKE~\cite{BonehBoyen08})~; l'opposabilit\'e, elle, d\'epend de l'interpr\'etation
institutionnelle et n'est pas minor\'ee par ces r\'eductions.

\subsubsection{Architecture de la couche Clay et budget entropique}
\label{subsubsec:couche_clay}

La couche Clay encode le contexte $c$ sous une forme maximisant l'information extractible
par l'interpr\`ete institutionnel tout en maintenant les garanties Gold et Iron. Chacun des
trois r\'egimes appelle une technique propre.

Pour le r\'egime rare ($H(\mathcal{C}) < H^*(\lambda)$), l'encodage d\'eterministe
$\phi:\mathcal{C} \to \{0,1\}^{H(\mathcal{C})}$ encode l'int\'egralit\'e du contexte sans
perte~\cite{Shannon48}. Pour la zone de transition, la compression Johnson-Lindenstrauss~\cite{JohnsonLindenstrauss84,Achlioptas2003} projette le
contexte depuis la dimension $d = \Theta(\Delta \cdot H(\mathcal{C}))$ vers $k =
O(\log n/\epsilon^2)$ dimensions, pr\'eservant les distances \`a $\epsilon$ pr\`es. Pour le
r\'egime riche enfin, le hachage Poseidon~\cite{Poseidon2019,Grassi2021} compresse un contexte arbitrairement complexe en
un condens\'e de taille fixe ($256$ bits), agr\'ement\'e de m\'etadonn\'ees structur\'ees pour les
recherches institutionnelles.

Le budget entropique total de chaque preuve SOTP s'\'ecrit
\begin{equation}
H_{\text{total}} = H_{\text{Gold}} + H_{\text{Iron}} + H_{\text{Clay}}
  \approx 2\lambda + H(\mathcal{C})
\label{eq:entropy_budget}
\end{equation}
Pour $\lambda = 128$ et $H(\mathcal{C}) = 40$ bits (vote type Helios~\cite{Adida08}), cela donne
$\approx 296$ bits. Ce budget contraint la couche Clay, car pr\'eserver trop de contexte
peut corr\'eler $\Sigma$ avec $m$, r\'eduisant $\Conf$~; compresser trop agressivement
r\'eduit $\Opp$ parce que l'interpr\`ete ne peut plus extraire les verdicts~\cite{Smith09,Alvim12}.

L'allocation \'equilibr\'ee est modulaire par conception~\cite{canetti2001,Goldreich02}. La couche Gold peut substituer
AES-GCM~\cite{AES01} \`a ChaCha20~\cite{ChaCha20} sans affecter Iron ou Clay~; la couche Iron peut passer d'un
commitment Pedersen \`a un accumulateur RSA~\cite{CamenischLysyanskaya02} sans toucher Gold ou Clay~; la couche Clay
peut changer de technique d'encodage ind\'ependamment. Cette modularit\'e facilite les
migrations quantiques et l'adaptation aux \'evolutions r\'eglementaires.
\end{sloppypar}

% -----------------------------------------------------------------------
\subsection{SOTP-Low~: r\'egime \`a contexte rare}
\label{subsec:sotp_low}
% -----------------------------------------------------------------------
\begin{sloppypar}
Nous proposons, pour le r\'egime \`a contexte rare, une construction
d\'eterministe simple dont nous \'etablissons les propri\'et\'es cryptographiques
(Remarque~\ref{rem:statut_sotp}).

\subsubsection{Motivation et champ d'application}
\label{subsubsec:motivation_low}

Le r\'egime cible est celui des contextes que l'on peut encoder int\'egralement dans
la preuve, $H(\mathcal{C}) < H^*(\lambda)$, avec $H^*(\lambda) = \sqrt{\lambda}$ comme
crit\`ere de construction (environ $11$ bits \`a $\lambda = 128$), distinct du seuil
$H^\sharp$ du th\'eor\`eme (Remarque~\ref{rem:activite_delta}). Les signatures contractuelles ECDSA~\cite{Johnson01,NIST.FIPS.186-5} en sont le cas d'\'ecole, puisqu'un contrat
standardis\'e et ses parties sont souvent param\'etr\'es par un contexte d'entropie
$H(\mathcal{C}) \approx 20$ bits~\cite{Mason07,Kerr04}. D\`es lors que $\Delta \gtrsim H(\mathcal{C})$, la
capacit\'e institutionnelle absorbe l'int\'egralit\'e du contexte de sorte que l'opposabilit\'e peut
approcher~1.

L'id\'ee consiste \`a encoder le contexte \emph{d\'eterministiquement} dans la couche Clay,
ce qui n'entra\^ine ni perte de compression ni circuit co\^uteux. La preuve se compose d'un masque PRG
(Gold), d'un commitment Pedersen (Iron) et d'un triplet $(\phi(c), \mathsf{Hash}(c), V(c))$
(Clay), cette derni\`ere composante r\'ev\`ele le contexte sans r\'ev\'eler le message, puisque
$M \perp C$ par construction.

\subsubsection{Construction SOTP-Low}
\label{subsubsec:construction_low}

\begin{algorithm}[H]
\caption{SOTP-Low}
\label{alg:sotp_low}
\begin{algorithmic}[1]
\State \textbf{Param\`etres} : $\lambda$, $H(\mathcal{C}) < H^*(\lambda)$, $(\Delta, \alpha, \beta)$
\State \textbf{$\mathsf{Setup}(1^\lambda)$} :
\State \quad G\'en\'erer groupe cyclique $(\mathbb{G}, q)$ avec $|q| \geq 2\lambda$ et DLog difficile
\State \quad \'echantillonner $g, h \xleftarrow{\$} \mathbb{G}$ tels que $\log_g h$ inconnu
\State \quad Choisir PRG $\mathsf{PRG}: \mathbb{Z}_q \to \{0,1\}^\lambda$
\State \quad Choisir $\mathsf{Hash}: \mathcal{C} \to \{0,1\}^\lambda$ r\'esistante aux collisions
\State \quad D\'efinir encodage d\'eterministe $\phi: \mathcal{C} \to \{0,1\}^{H(\mathcal{C})}$
\State \quad \textbf{Retourner} $pk = (\mathbb{G}, q, g, h, \mathsf{PRG}, \mathsf{Hash}, \phi)$
\State \textbf{$\mathsf{Prove}(pk, m, c)$} :
\State \quad \textcolor{blue}{// Gold}
\State \quad $s \xleftarrow{\$} \mathbb{Z}_q$ ;\; $\sigma_1 \leftarrow m \oplus \mathsf{PRG}(s)$
\State \quad \textcolor{blue}{// Iron}
\State \quad $\rho \xleftarrow{\$} \mathbb{Z}_q$ ;\; $\sigma_2 \leftarrow g^{s} \cdot h^\rho$~\cite{Pedersen91}
\State \quad \textcolor{blue}{// Clay}
\State \quad $\sigma_3 \leftarrow (\phi(c),\; \mathsf{Hash}(c),\; V(c))$
\State \quad \textbf{Retourner} $\sigma = (\sigma_1, \sigma_2, \sigma_3)$ ;\; t\'emoin secret $(s, \rho)$
\State \textbf{$\mathsf{Verify}(pk, \sigma, c)$} :
\State \quad V\'erifier $\mathsf{Hash}(c) \stackrel{?}{=} \sigma_3.\mathsf{hc}$ et
         $\phi(c) \stackrel{?}{=} \sigma_3.h_c$ et $\sigma_2 \in \mathbb{G}$
\State \textbf{$\mathsf{Open}(pk, \sigma, c, (m, s, \rho))$} :
\State \quad Accepter ssi $\sigma_2 = g^{s} h^{\rho}$ et $\sigma_1 = m \oplus \mathsf{PRG}(s)$
\end{algorithmic}
\end{algorithm}

L'algorithme op\`ere en trois phases distinctes qui correspondent exactement
aux trois couches du paradigme Gold-Iron-Clay.

La phase \textsf{Setup} installe les param\`etres cryptographiques partag\'es.
Elle g\'en\`ere un groupe cyclique $(\mathbb{G}, q)$ d'ordre premier $\lambda$
bits dans lequel le logarithme discret est difficile, ce qui fonde
la s\'ecurit\'e de la couche Iron. Elle tire al\'eatoirement deux g\'en\'erateurs
$g$ et $h$ dont le rapport discret $\log_g h$ est inconnu de tous les
participants~; cette inconnaissance rend l'engagement de Pedersen
statistiquement cachant. Elle fixe enfin trois primitives fonctionnelles, \`a savoir
un PRG pour le masquage de message, une fonction de hachage r\'esistante
aux collisions pour l'authentification du contexte et un encodage
d\'eterministe $\phi$ qui serialise le contexte institutionnel $c$ sur
exactement $H(\mathcal{C})$ bits. Ce dernier point est sp\'ecifique au
r\'egime \`a contexte rare. En effet, parce que $H(\mathcal{C}) < H^*(\lambda)$, la
totalit\'e du contexte tient dans cet encodage sans compression avec perte.

La phase \textsf{Prove} construit la preuve en trois \'etapes parall\`eles
ind\'ependantes. La couche Gold (ligne~11) tire une graine
secr\`ete $s$ uniform\'ement dans $\mathbb{Z}_q$ et produit le message
masqu\'e $\sigma_1 = m \oplus \mathsf{PRG}(s)$. Le message r\'eel est ainsi
int\'egralement dissimul\'e derri\`ere un pad pseudoal\'eatoire, rendant
$\sigma_1$ indistinguable d'une cha\^ine uniforme pour tout adversaire ne
connaissant pas $s$. La couche Iron (ligne~13) construit
un engagement de Pedersen sur la graine $s$, soit $\sigma_2 = g^{s} \cdot h^\rho$
o\`u $\rho$ est un scalaire al\'eatoire. Tirer la graine dans $\mathbb{Z}_q$ est
n\'ecessaire. Avec $s \in \{0,1\}^\lambda$ et un engagement sur $\mathsf{int}(s)
\bmod q$, deux graines $s$ et $s+q$ donneraient en effet le m\^eme engagement, ce qui ferait perdre le
caract\`ere liant. Cet
engagement lie le producteur de la preuve \`a la valeur de $s$ sans la
r\'ev\'eler~; il garantit que le message masqu\'e $\sigma_1$ ne peut \^etre
ouvert que sur un unique message $m$. La couche Clay
(ligne~15) produit le triplet $\sigma_3 = (\phi(c), \mathsf{Hash}(c), V(c))$.
L'encodage d\'eterministe $\phi(c)$ y transmet la structure s\'emantique du
contexte, le condens\'e $\mathsf{Hash}(c)$ en garantit l'int\'egrit\'e tandis que
$V(c)$ est le verdict institutionnel attendu. Cette couche porte
l'int\'egralit\'e de l'information contextuelle exploitable par l'arbitre.
La preuve compl\`ete $\sigma = (\sigma_1, \sigma_2, \sigma_3)$ est
retourn\'ee publiquement~; la paire $(s, \rho)$ reste secr\`ete chez le
producteur.

La phase \textsf{Verify} est d\'elib\'er\'ement l\'eg\`ere, puisqu'elle contr\^ole
uniquement que le condens\'e de hachage du contexte soumis correspond \`a
celui encod\'e dans $\sigma_3$, que l'encodage $\phi(c)$ co\"incide et que
$\sigma_2$ appartient bien au groupe $\mathbb{G}$. Elle ne v\'erifie pas la
coh\'erence entre $\sigma_1$ et $\sigma_2$, car ce lien n'est contr\^ol\'e que lors
de l'ouverture ($\mathsf{Open}$), qui requiert le t\'emoin $(m, s, \rho)$~;
sans ouverture, $\sigma_1$ et $\sigma_2$ ne sont pas li\'es, de sorte que
$\mathsf{Verify}$ seul accepte toute paire bien form\'ee. L'arbitre peut
donc v\'erifier la forme de la preuve sans acc\'eder au message $m$~; la
garantie de liaison porte sur l'ouverture. La couche Clay pr\'esente le
contexte sous une forme lisible, mais, \'etant une fonction du contexte public
$c$, elle n'apporte aucune information conditionnelle au sens de $\Opp$
(Remarque~\ref{rem:statut_sotp}).

\subsubsection{Propri\'et\'es de SOTP-Low}
\label{subsubsec:saturation_low}

\begin{proposition}[Propri\'et\'es cryptographiques de SOTP-Low]\label{thm:saturation_sotp_low}
Sous s\'ecurit\'e PRG~\cite{BlumMicali84,Yao82} et difficult\'e DLog~\cite{DiffieHellman76}~:
\begin{enumerate}
\item tout adversaire PPT a un avantage IND-CONF n\'egligeable contre SOTP-Low
      (confidentialit\'e calculatoire)~;
\item la compl\'etude est parfaite et, pour la proc\'edure d'ouverture,
      la probabilit\'e de gain BIND-REL est n\'egligeable, donc
      $\Relver(\mathrm{SOTP\text{-}Low}) = 1 - \negl(\lambda)$ au sens de la
      Remarque~\ref{rem:jeu_bind_rel_link}.
\end{enumerate}
\end{proposition}

\begin{remark}[Heuristique de conception de SOTP-Low (non prouv\'ee)]
\label{rem:facteur_limitant_low}
La formulation initiale de ce r\'esultat \'enon\c{c}ait la minoration
$\Conf \cdot \Rel \cdot \Opp \geq 0{,}75 \cdot \min\{\Delta, H(\mathcal{C})\}/\lambda$.
La formalisation a conduit \`a la remplacer, pour trois raisons. (i)~La composante $\sigma_3 = (\phi(c), \mathsf{Hash}(c), V(c))$
est une fonction du contexte public, donc $I(V;\sigma_3|C) = 0$~; (ii)~le passage
par $H(V|C) \approx \lambda$ contredit l'Hypoth\`ese~\ref{ass:verdicts_efficace}
($\log_2|\mathcal{V}_\lambda| \leq \alpha H(\mathcal{C})$)~; (iii)~une telle minoration
d\'epasserait le majorant du Th\'eor\`eme~\ref{thm:trilemme_cro_conditionnel} d\`es
que $\gamma > 4/3$ tout en portant sur $\Relver$ au lieu de $\Relrec$. L'expression
$0{,}75\,\Delta/\lambda$ est conserv\'ee comme \emph{heuristique de conception}, car
elle d\'ecrit la fraction de la capacit\'e d'extraction qu'un encodage entier de
$\lfloor H(\mathcal{C}) \rfloor$ bits permettrait d'exploiter si ce contexte
\'etait un contexte d'esp\`ece inconnu de l'institution ($1 - 1/\Delta \geq 0{,}75$
pour $\Delta \geq 4$). Sa justification rigoureuse, dans un mod\`ele \`a contexte
priv\'e, est une perspective (chapitre~\ref{chap:resultats_discussions}).
\end{remark}

\begin{Preuve}[Proposition~\ref{thm:saturation_sotp_low}~: propri\'et\'es cryptographiques de SOTP-Low]\label{prv:saturation_sotp_low}

\begin{cmtEtape}{Vue d'ensemble.}
SOTP-Low est constitu\'e de trois couches, Gold (masquage par PRG),
Iron (engagement de Pedersen sur la graine $s$) et Clay (encodage
d\'eterministe $\phi(c)$). La confidentialit\'e et la fiabilit\'e de
v\'erification sont \'etablies s\'epar\'ement.
\end{cmtEtape}

\noindent\textbf{Rappel de la construction.}
\begin{align}
\sigma_1 &= m \oplus \mathsf{PRG}(s),
\quad s \xleftarrow{\$} \mathbb{Z}_q
\tag{couche Gold}\\
\sigma_2 &= g^{s} \cdot h^\rho,
\quad \rho \xleftarrow{\$} \mathbb{Z}_q
\tag{couche Iron}\\
\sigma_3 &= \bigl(\phi(c),\; \mathsf{Hash}(c),\; V(c)\bigr)
\tag{couche Clay}
\end{align}
o\`u $(g, h)$ sont des g\'en\'erateurs du groupe $(\mathbb{G}, q)$ avec
$\log_g h$ inconnu. Le t\'emoin secret est $(s, \rho)$.

\medskip
\noindent\textbf{Partie I~. Confidentialit\'e calculatoire.}
La loi de $\sigma$ sachant $M = m_i$ ne d\'epend de $m_i$ que par
$\sigma_1$. En effet, $\sigma_2$ est parfaitement masquant (pour tout $s$, $g^s h^\rho$
est uniforme dans $\mathbb{G}$ lorsque $\rho$ est uniforme) tandis que $\sigma_3$ ne
d\'epend que de $c$, ind\'ependant de $m$ ($M \perp C$).
\cmt{Seule $\sigma_1$ porte de l'information sur $m$.}

Supposons qu'un distingueur PPT $\mathcal{D}$ distingue $m_0 \oplus
\mathsf{PRG}(s)$ de $m_1 \oplus \mathsf{PRG}(s)$ avec un avantage non
n\'egligeable $\mu(\lambda)$. L'adversaire $\mathcal{B}$ qui re\c{c}oit
$y$ (sortie du PRG ou cha\^ine uniforme), tire $b$, soumet $m_b \oplus y$ \`a
$\mathcal{D}$ et r\'epond \og PRG\fg{} si $\mathcal{D}$ retrouve $b$, distingue le
PRG avec un avantage $\mu(\lambda)/2$, puisque si $y$ est uniforme, $m_b \oplus y$ ne
contient aucune information sur $b$. Cela contredit la s\'ecurit\'e du
PRG~\cite{BlumMicali84}. L'avantage IND-CONF de tout adversaire PPT est donc
n\'egligeable.
\cmt{La confidentialit\'e est calculatoire. Au sens statistique de la D\'efinition~\ref{def:confidentialite_clo}, elle n'est garantie que si le PRG est remplac\'e par un masque uniforme.}

\medskip
\noindent\textbf{Partie II~. Fiabilit\'e de v\'erification.}

\textit{Compl\'etude.} Pour toute preuve honn\^ete, $\mathsf{Hash}(c) = \sigma_3.\mathsf{hc}$
et $\phi(c) = \sigma_3.h_c$ par construction, $\sigma_2 \in \mathbb{G}$ par
fermeture du groupe~; enfin $\mathsf{Open}$ accepte le t\'emoin honn\^ete $(m, s, \rho)$.
La compl\'etude est parfaite.
\cmtFormule{La v\'erification et l'ouverture acceptent toute preuve honn\^ete.}

\textit{Liaison (r\'eduction \`a DLog).} Supposons qu'un adversaire PPT
produise une preuve $\sigma$ et deux ouvertures accept\'ees
$(m_0, s_0, \rho_0)$ et $(m_1, s_1, \rho_1)$ avec $m_0 \neq m_1$. Comme
$\sigma_1 = m_0 \oplus \mathsf{PRG}(s_0) = m_1 \oplus \mathsf{PRG}(s_1)$ et
$m_0 \neq m_1$, on a $s_0 \neq s_1$ dans $\mathbb{Z}_q$. L'\'egalit\'e
$g^{s_0} h^{\rho_0} = g^{s_1} h^{\rho_1}$ impose alors $\rho_0 \neq \rho_1$ et
\begin{equation}
\log_g h = \frac{s_0 - s_1}{\rho_1 - \rho_0} \bmod q,
\end{equation}
\cmtFormule{Deux ouvertures distinctes r\'ev\`elent $\log_g h$.}
ce qui contredit la difficult\'e du logarithme discret~\cite{DiffieHellman76}.
La probabilit\'e de gain BIND-REL est donc n\'egligeable, d'o\`u
$\Relver(\mathrm{SOTP\text{-}Low}) = 1 - \negl(\lambda)$.
\cmt{C'est ici que le tirage de $s$ dans $\mathbb{Z}_q$ est indispensable.}
\end{Preuve}

\subsubsection{Support de validation~: signatures ECDSA}
\label{subsubsec:validation_ecdsa}

\`A titre d'illustration, sans valeur de validation (les m\'etriques CLO d'ECDSA sont
des estimations param\'etriques analytiques issues des sp\'ecifications
publi\'ees~\cite{Johnson01,NIST.FIPS.186-5} et ECDSA n'est pas une instance du
th\'eor\`eme, Remarque~\ref{rem:portee_cas_etude}), les param\`etres $\Delta = 35$ bits,
$\lambda = 128$ donnent pour l'heuristique de SOTP-Low la valeur
$0{,}75 \times 35/128 \approx 0{,}205$. Cette valeur n'est pas comparable au produit
CLO d'ECDSA ($0{,}273$).
\end{sloppypar}

% -----------------------------------------------------------------------
\subsection{SOTP-Base~: r\'egime de transition}
\label{subsec:sotp_base}
% -----------------------------------------------------------------------
\begin{sloppypar}
Pour les contextes interm\'ediaires, nous proposons une construction fond\'ee sur la
compression Johnson-Lindenstrauss~\cite{JohnsonLindenstrauss84}, qui r\'eduit la taille de la
repr\'esentation du contexte tout en pr\'eservant ses distances \`a $\varepsilon$ pr\`es.

\subsubsection{Motivation et champ d'application}
\label{subsubsec:motivation_base}

Helios~\cite{Adida08,Adida09UCL} est l'ach\'etype m\^eme de ce r\'egime, puisque le contexte \'electoral (registre des \'electeurs, r\`egles du
scrutin, proc\`es-verbal) repr\'esente $H(\mathcal{C}) \approx 40$--$50$ bits pour une
\'election de taille moyenne, la capacit\'e d'un comit\'e de garants sp\'ecialis\'es \'etant estim\'ee
\`a $\Delta \approx 56$ bits. L'encodage d\'eterministe de SOTP-Low ne passe pas \`a
l'\'echelle, car la preuve serait trop volumineuse. Les zk-SNARKs de SOTP-High
seraient \`a l'inverse surdimensionn\'es pour ce niveau d'entropie.

La compression Johnson-Lindenstrauss~\cite{JohnsonLindenstrauss84,DasguptaGupta03,Achlioptas2003} permet de
projeter le contexte depuis la dimension $d = \Theta(\Delta \cdot H(\mathcal{C}))$ vers
$k = O(\log n/\epsilon^2)$ dimensions via une matrice al\'eatoire $\Phi$, pr\'eservant les
distances \`a un facteur $1 \pm \epsilon$ pr\`es. Que l'information extractible soit
pr\'eserv\'ee dans la m\^eme proportion rel\`eve d'une hypoth\`ese de conception que le
lemme ne garantit pas.

\subsubsection{Construction SOTP-Base}
\label{subsubsec:construction_base}

\begin{algorithm}[H]
\caption{SOTP-Base~: construction avec compression JL}
\label{alg:sotp_base}
\begin{algorithmic}[1]
\State \textbf{Param\`etres} : $\lambda$, $H^*(\lambda) \leq H(\mathcal{C}) \leq \lambda/\log\lambda$, $\varepsilon = 0{,}1$
\State \textbf{$\mathsf{Setup}(1^\lambda)$} :
\State \quad G\'en\'erer $(\mathbb{G}, q, g, h)$, PRG, Hash comme SOTP-Low
\State \quad Fixer $d \leftarrow \lceil 4\Delta H(\mathcal{C})/\varepsilon^2 \rceil$
\State \quad Fixer $k \leftarrow \lceil 8\log(2^{H(\mathcal{C})})/\varepsilon^2 \rceil$
\State \quad G\'en\'erer matrice JL $\Phi \xleftarrow{\$} \mathcal{N}(0, 1/k)^{k \times d}$~\cite{DasguptaGupta03,Vempala05}
\State \quad \textbf{Retourner} $pk = (\mathbb{G}, q, g, h, \mathsf{PRG}, \mathsf{Hash}, \Phi, k, d)$
\State \textbf{$\mathsf{Prove}(pk, m, c)$} :
\State \quad Couches Gold et Iron identiques \`a SOTP-Low
\State \quad \textcolor{blue}{// Clay JL}
\State \quad $\vec{c} \leftarrow \mathsf{embed}(c) \in \mathbb{R}^d$
\State \quad $\sigma_3 \leftarrow \Phi \vec{c} \in \mathbb{R}^k$
\State \quad \textbf{Retourner} $\sigma = (\sigma_1, \sigma_2, \sigma_3)$
\State \textbf{$\mathsf{Verify}(pk, \sigma, c)$} :
\State \quad V\'erifier $\|\Phi \cdot \mathsf{embed}(c) - \sigma_3\| \leq \varepsilon$
\State \quad V\'erifier $\sigma_2 \in \mathbb{G}$ et coh\'erence Gold
\end{algorithmic}
\end{algorithm}

SOTP-Base partage avec SOTP-Low les couches Gold et Iron inchang\'ees.
La diff\'erence enti\`ere r\'eside dans la couche Clay, qui remplace
l'encodage d\'eterministe bijectif par une compression
Johnson-Lindenstrauss. Cette substitution est rendue n\'ecessaire par
le r\'egime contextuel. Ici $H(\mathcal{C}) \geq H^*(\lambda)$, ce
qui signifie que le contexte institutionnel est trop riche pour tenir
enti\`erement dans un encodage de taille fixe sans perte.

La phase \textsf{Setup} h\'erite du groupe, du PRG et du Hash de
SOTP-Low, puis calcule deux dimensions caract\'eristiques de la
compression. La dimension cible $d = \lceil 4\Delta
H(\mathcal{C})/\varepsilon^2 \rceil$ (ligne~4) calibre l'espace de
repr\'esentation interm\'ediaire sur le produit de la capacit\'e
d'extraction institutionnelle $\Delta$ et de l'entropie contextuelle, car
plus le contexte est riche et plus l'institution peut en extraire,
plus la repr\'esentation doit \^etre large pour pr\'eserver les
distances s\'emantiques pertinentes. La dimension de sortie
$k = \lceil 8\log(2^{H(\mathcal{C})})/\varepsilon^2 \rceil$
(ligne~5) d\'etermine la taille finale de la couche Clay transmise
dans la preuve~; elle est calibr\'ee pour que le lemme JL garantisse
la pr\'eservation des distances \`a $\varepsilon = 0{,}1$ pr\`es avec
haute probabilit\'e. La matrice $\Phi$ (ligne~6) est tir\'ee
al\'eatoirement selon une gaussienne centr\'ee de variance $1/k$. 
Cette normalisation assure que les normes euclidiennes sont pr\'eserv\'ees
en esp\'erance, propri\'et\'e sur laquelle repose la
garantie JL.

La phase \textsf{Prove} (lignes~8--13) reconstruit $\sigma_1$ et
$\sigma_2$ exactement comme SOTP-Low, puis traite le contexte
diff\'eremment. La ligne~11 plonge d'abord le contexte $c$ dans
$\mathbb{R}^d$ via la fonction $\mathsf{embed}$, qui traduit la
structure s\'emantique du contexte institutionnel en vecteur de r\'eels, de sorte que
deux contextes proches s\'emantiquement produisent des vecteurs proches
au sens de la distance euclidienne. La ligne~12 applique ensuite la
compression $\sigma_3 = \Phi\vec{c} \in \mathbb{R}^k$, par laquelle le vecteur de
$d$ dimensions est projet\'e vers un espace de $k \ll d$ dimensions.
Cette projection constitue la couche Clay de SOTP-Base. La
preuve transmise $\sigma_3$ est donc la projection compress\'ee du contexte,
distincte du contexte lui-m\^eme et de son encodage exact. Les distances
s\'emantiques entre contextes distincts y sont pr\'eserv\'ees \`a
$\varepsilon$ pr\`es par la garantie JL. La pr\'eservation de l'information
extractible reste une hypoth\`ese de conception.

La phase \textsf{Verify} (lignes~14--16) refl\`ete cette logique, puisqu'elle v\'erifie la proximit\'e $\|\Phi \cdot
\mathsf{embed}(c) - \sigma_3\| \leq \varepsilon$ au lieu de l'\'egalit\'e exacte
$\Phi \cdot \mathsf{embed}(c) = \sigma_3$. Ce crit\`ere souple
est in\'evitable, car la compression avec perte $\varepsilon = 0{,}1$
signifie que le v\'erificateur ne peut retrouver $\sigma_3$ qu'\`a
$10\%$ pr\`es depuis le contexte re\c{c}u. Le choix de $\varepsilon$
est ainsi un param\`etre de conception explicite qui gouverne le
compromis entre la fid\'elit\'e de la repr\'esentation contextuelle
et la dimension $k$ de la preuve transmise. R\'eduire $\varepsilon$
augmente la fid\'elit\'e mais exige un $k$ plus grand, alourdissant
la preuve. Avec $\varepsilon = 0{,}1$, $\Delta = 56$ et $H(\mathcal{C}) = 40$, la
dimension interm\'ediaire vaut $d = \lceil 4\Delta H(\mathcal{C})/\varepsilon^2 \rceil
= 896\,000$.

\subsubsection{Propri\'et\'es de SOTP-Base}
\label{subsubsec:saturation_base}

\begin{proposition}[Propri\'et\'es cryptographiques de SOTP-Base]\label{thm:saturation_sotp_base}
Sous s\'ecurit\'e PRG~\cite{BlumMicali84} et difficult\'e DLog~\cite{DiffieHellman76}, les
couches Gold et Iron de SOTP-Base \'etant identiques \`a celles de SOTP-Low,
SOTP-Base v\'erifie les deux propri\'et\'es de la
Proposition~\ref{thm:saturation_sotp_low}, \`a savoir la confidentialit\'e calculatoire et
$\Relver = 1 - \negl(\lambda)$ pour la proc\'edure d'ouverture. De plus, pour
tous contextes $c, c'$, avec probabilit\'e au moins $1 - \delta$ sur le tirage
de $\Phi$, $\|\Phi\,\mathsf{embed}(c) - \Phi\,\mathsf{embed}(c')\|^2 \in
(1 \pm \varepsilon)\|\mathsf{embed}(c) - \mathsf{embed}(c')\|^2$
(Lemme~\ref{lem:johnson_lindenstrauss}).
\end{proposition}

\begin{proof}
Les arguments de la Proposition~\ref{thm:saturation_sotp_low} s'appliquent sans
changement aux couches Gold et Iron~; la couche Clay n'intervient ni dans la
confidentialit\'e (elle ne d\'epend que de $c$, ind\'ependant de $m$) ni dans la
liaison. La derni\`ere propri\'et\'e est le lemme de Johnson--Lindenstrauss appliqu\'e
\`a l'ensemble fini des plongements de contextes.
\end{proof}

\begin{remark}[Heuristique de conception de SOTP-Base (non prouv\'ee)]
\label{rem:heuristique_sotp_base}
La formulation initiale de ce r\'esultat \'enon\c{c}ait
$\Conf \cdot \Rel \cdot \Opp \geq 0{,}5 \cdot \min\{\Delta/d,\, \beta\log_2 d/d\}$.
La formalisation a conduit \`a la remplacer pour les raisons expos\'ees \`a la
Remarque~\ref{rem:facteur_limitant_low}. En effet, $\sigma_3 = \Phi\,\mathsf{embed}(c)$ est
une fonction du contexte public, donc $I(V;\sigma_3|C) = 0$~; la pr\'eservation
des distances n'implique pas celle de l'information extractible~; enfin une
minoration ne peut se d\'eduire d'une majoration. L'expression
$0{,}5 \cdot \min\{\Delta/d, \Delta/(\beta\log_2 d)\}$ est conserv\'ee comme
heuristique de dimensionnement de la couche Clay.
\end{remark}

\subsubsection{Support de validation~: vote \'electronique Helios}
\label{subsubsec:validation_helios}

\`A titre d'illustration, l'heuristique de SOTP-Base pour les param\`etres
analytiques de Helios~\cite{Adida08,MINKA2026CRO} ($\Delta = 56$ bits, $\beta = 1{,}4$,
$d = 560$) donne $0{,}5 \times \min\{56/560,\, 56/(1{,}4 \log_2 560)\} \approx 0{,}050$.
Cette valeur n'est pas comparable au produit CLO de Helios ($0{,}437$), qui porte
sur $\Relver$ et n'est pas une instance du th\'eor\`eme
(Remarque~\ref{rem:portee_cas_etude}). La dimension $d = 560$ retenue ici est un
choix de dimensionnement, distinct de la r\`egle $d = \lceil 4\Delta H(\mathcal{C})/\varepsilon^2\rceil$.
\end{sloppypar}

% -----------------------------------------------------------------------
\subsection{SOTP-High~: r\'egime \`a contexte riche}
\label{subsec:sotp_high}
% -----------------------------------------------------------------------
\begin{sloppypar}
Pour les contextes riches, nous proposons une construction fond\'ee sur les zk-SNARKs
succincts Groth16, qui prot\`ege le t\'emoin (divulgation nulle de connaissance) au prix
d'une opposabilit\'e r\'eduite.

\subsubsection{Motivation et champ d'application}
\label{subsubsec:motivation_high}

La conformit\'e financi\`ere via zk-SNARKs~\cite{Groth16,Garratt15,BaselIII} illustre ce r\'egime, o\`u transactions complexes,
r\'eglementations multi-juridictionnelles et historiques longs produisent $H(\mathcal{C})
\approx 70$--$90$ bits. M\^eme la compression JL de SOTP-Base serait insuffisante, car la
dimension cible $k = \Omega(H(\mathcal{C}))$ resterait trop large pour absorption par
un interpr\`ete \`a capacit\'e $\Delta$ born\'ee~\cite{Simon82,Miller56}.

Les zk-SNARKs succincts~(Groth16~\cite{Groth16,BCTV14,Gennaro13}) l\`event cet obstacle en produisant une preuve
de taille constante ($128$ \`a $192$ octets selon la courbe) ind\'ependante de la complexit\'e du
circuit v\'erifi\'e. Le caract\`ere succinct brise le couplage lin\'eaire entre $H(\mathcal{C})$ et la
taille de preuve, permettant d'encoder des relations complexes sur un contexte de
plusieurs milliers de bits dans un condens\'e que l'interpr\`ete peut v\'erifier en quelques millisecondes.

Cette solution a une double contrepartie. Une \'etape de mise en place de confiance (\emph{trusted
setup}) est requise~\cite{Bowe2017}. L'opposabilit\'e est aussi r\'eduite, plafonn\'ee par la perte
entropique du hachage compressant le contexte en 256 bits.

\subsubsection{Construction SOTP-High}
\label{subsubsec:construction_high}

\begin{algorithm}[H]
\caption{SOTP-High~: construction avec zk-SNARKs Groth16}
\label{alg:sotp_high}
\begin{algorithmic}[1]
\State \textbf{Param\`etres} : $\lambda$, $H(\mathcal{C}) > \lambda/\log\lambda$, circuit $\mathcal{C}_{\text{circuit}}$
\State \textbf{$\mathsf{Setup}(1^\lambda, \mathcal{C}_{\text{circuit}})$} :
\State \quad \'echantillonner $\tau, \alpha_s, \beta_s, \gamma, \delta \xleftarrow{\$} \mathbb{F}_p^*$ (toxiques)
\State \quad Construire $pk_{\mathsf{prove}}$ et $vk_{\mathsf{verify}}$ (CRS Groth16)~\cite{Groth16,BCTV14}
\State \quad \textcolor{red}{// D\'etruire toxiques~: $\tau, \alpha_s, \beta_s, \gamma, \delta$}
\State \quad $(g, h) \xleftarrow{\$} \mathbb{G}^2$ pour couche Iron
\State \quad \textbf{Retourner} $pk = (pk_{\mathsf{prove}}, vk_{\mathsf{verify}}, g, h)$
\State \textbf{$\mathsf{Prove}(pk, m, c)$} :
\State \quad \textcolor{blue}{// Gold + Iron (Groth16 + Pedersen)}
\State \quad $\rho \xleftarrow{\$} \mathbb{Z}_q$ ;\; $c_m \leftarrow g^m \cdot h^\rho$~\cite{Pedersen91}
\State \quad $h_c \leftarrow \mathsf{Poseidon}(c)$~\cite{Poseidon2019,Grassi2021} ;\; $\pi \leftarrow \mathsf{Groth16.Prove}(pk_{\mathsf{prove}}, (h_c, c_m), (m, c, \rho))$
\State \quad \textit{relation~: $h_c = \mathsf{Poseidon}(c) \wedge c_m = g^m h^\rho \wedge \mathsf{Conforme}(m, c) = 1$}
\State \quad \textcolor{blue}{// Clay~: condens\'e hash + m\'etadonn\'ees}
\State \quad \textbf{Retourner} $\sigma = (h_c, \pi, c_m, \mathsf{metadata})$
\State \textbf{$\mathsf{Verify}(\sigma)$} :
\State \quad $b_1 \leftarrow \mathsf{Groth16.Verify}(vk_{\mathsf{verify}}, (h_c, c_m), \pi)$
\State \quad $b_2 \leftarrow [c_m \in \mathbb{G}]$
\State \quad \textbf{Retourner} $b_1 \wedge b_2$
\end{algorithmic}
\end{algorithm}

SOTP-High est la construction la plus complexe des trois. Elle est
con\c{c}ue pour le r\'egime o\`u le contexte institutionnel est
arbitrairement riche, $H(\mathcal{C}) > \lambda/\log\lambda$, ce qui
rend impossible toute repr\'esentation fid\`ele du contexte dans la
preuve sans r\'ev\'eler de l'information confidentielle. La solution
adopt\'ee s'\'ecarte donc de SOTP-Low et SOTP-Base. Au lieu de transmettre le contexte ou sa compression, on transmet une
\emph{preuve de connaissance} du contexte via un zk-SNARK Groth16.

La phase \textsf{Setup} (lignes~2--7) est la plus sensible des trois
constructions. Elle \'echantillonne cinq scalaires $\tau, \alpha_s,
\beta_s, \gamma, \delta$ dans le corps premier $\mathbb{F}_p^*$, qui
sont les param\`etres toxiques de l'\'etape de configuration
Groth16. Ils servent \`a construire la cl\'e de preuve $pk_{\mathsf{prove}}$
et la cl\'e de v\'erification $vk_{\mathsf{verify}}$ qui forment le
Common Reference String (CRS). La ligne~5, signal\'ee en rouge, est
un imp\'eratif de s\'ecurit\'e. Si l'un quelconque de ces cinq
scalaires subsistait apr\`es le \textsf{Setup}, son d\'etenteur pourrait
forger des preuves valides pour n'importe quelle assertion, brisant
enti\`erement la propri\'et\'e de soundness du syst\`eme. En pratique,
cette destruction est r\'ealis\'ee via une \'etape multi-partite
o\`u la compromission simultan\'ee de tous les participants serait
n\'ecessaire pour reconstituer les toxiques. Le Setup produit enfin
une paire $(g, h)$ pour la couche Iron de Pedersen, ind\'ependante
du CRS Groth16.

La phase \textsf{Prove} (lignes~8--14) fusionne les trois couches
en deux op\'erations principales. Les lignes~10 et~11 r\'ealisent les
trois couches. Le hash Poseidon $h_c =
\mathsf{Poseidon}(c)$ comprime le contexte arbitrairement riche $c$
vers un condens\'e de 256 bits qui constitue l'entr\'ee publique du
circuit~; l'engagement de Pedersen $c_m = g^m \cdot h^\rho$ joue le r\^ole de
la couche Iron, puisqu'il lie le prouveur au message $m$ sans le r\'ev\'eler. La
preuve $\pi$ atteste que le prouveur conna\^it $(m, c, \rho)$ tels que
$h_c = \mathsf{Poseidon}(c)$, $c_m = g^m h^\rho$ et que le pr\'edicat de
conformit\'e $\mathsf{Conforme}(m, c)$ est satisfait~; cette relation
lie l'engagement au contexte et \`a l'assertion prouv\'ee. Cette preuve ne
r\'ev\`ele rien sur $(m, c)$ au-del\`a de ce que les entr\'ees publiques
r\'ev\`elent d\'ej\`a~; la propri\'et\'e de divulgation nulle de connaissance de
Groth16 en garantit la confidentialit\'e. La preuve finale $\sigma = (h_c, \pi, c_m,
\mathsf{metadata})$ est ainsi constitu\'ee du condens\'e contextuel,
de la preuve zk-SNARK, de l'engagement sur le message et de
m\'etadonn\'ees permettant \`a l'arbitre de situer la preuve dans son
contexte institutionnel.

La phase \textsf{Verify} (lignes~15--18) est l\'eg\`ere
au regard de la richesse du contexte original. Elle r\'eduit \`a deux
v\'erifications bool\'eennes ind\'ependantes. La premi\`ere, $b_1$, v\'erifie que la
preuve Groth16 $\pi$ est valide pour le condens\'e public $h_c$ via
la cl\'e $vk_{\mathsf{verify}}$~; la seconde, $b_2$, v\'erifie que $c_m$ est un
point l\'egitime du groupe $\mathbb{G}$. La conjonction $b_1 \wedge
b_2$ suffit \`a \'etablir que le prouveur connaissait effectivement
un message $m$ et un contexte $c$ dont $\mathsf{Poseidon}(c) = h_c$,
sans que le v\'erificateur ait acc\`es \`a $m$ ou \`a
$c$. Le co\^ut du r\'egime \`a contexte riche se situe
l\`a. L'arbitre peut v\'erifier la coh\'erence de la
preuve, mais l'information contextuelle qu'il peut extraire de $h_c$
est born\'ee par la taille du condens\'e ($256$ bits). Elle est m\^eme nulle si $c$ est
d\'ej\`a connu de l'institution. Ce que la preuve apporte au verdict tient \`a
l'assertion certifi\'ee ($\mathsf{Conforme}(m,c) = 1$), non au condens\'e.

\subsubsection{Propri\'et\'es de SOTP-High}
\label{subsubsec:saturation_high}

\begin{proposition}[Propri\'et\'es cryptographiques de SOTP-High]\label{thm:saturation_sotp_high}
Sous les hypoth\`eses q-SDH et q-PKE (\emph{knowledge soundness} de
Groth16~\cite{BonehBoyen08,BCTV14}), la difficult\'e DLog et une c\'er\'emonie de mise en
place de confiance honn\^ete~\cite{Bowe2017}~:
\begin{enumerate}
\item la preuve $\pi$ est \`a divulgation nulle de connaissance parfaite et
      $c_m$ est parfaitement masquant, de sorte que $(\pi, c_m)$ ne r\'ev\`ele rien de $m$
      au-del\`a des entr\'ees publiques~;
\item la compl\'etude est parfaite et, par la \emph{knowledge soundness}, pour
      tout prouveur PPT qui fait accepter $\sigma$, il existe un
      \emph{extracteur} (qui rembobine le prouveur ou exploite la trappe de
      l'hypoth\`ese q-PKE, contrairement au v\'erificateur) produisant un t\'emoin
      $(m, c, \rho)$ satisfaisant la relation~; d'o\`u $\Relver = 1 - \negl(\lambda)$.
\end{enumerate}
\end{proposition}

\begin{remark}[Heuristique de conception de SOTP-High (non prouv\'ee)]
\label{rem:heuristique_sotp_high}
La formulation initiale de ce r\'esultat \'enon\c{c}ait
$\Conf \cdot \Rel \cdot \Opp \geq 0{,}45 \cdot \beta\log_2\lambda/\lambda$, avec une
constante $0{,}45$ \og calibr\'ee\fg{} sur les syst\`emes d\'eploy\'es. La formalisation a
conduit \`a la remplacer. Le condens\'e $h_c = \mathsf{Poseidon}(c)$ et les
m\'etadonn\'ees sont des fonctions du contexte public ($I(V;h_c|C) = 0$)~; trois
\'el\'ements de groupe repr\'esentent $\Theta(\lambda)$ bits et non $O(\log\lambda)$~;
le v\'erificateur ne peut extraire le t\'emoin, faute de quoi la preuve ne serait
pas \`a divulgation nulle~; enfin une constante ajust\'ee sur des donn\'ees n'a pas sa place
dans un \'enonc\'e de th\'eor\`eme. L'expression $0{,}45\,\beta\log_2\lambda/\lambda$ est
conserv\'ee comme heuristique de conception~; $\beta$ y d\'esigne la constante de
r\'egularit\'e du protocole (D\'efinition~\ref{def:beta_regularite}).
\end{remark}

\subsubsection{Support de validation~: conformit\'e zk-SNARKs}
\label{subsubsec:validation_snark}

Pour les param\`etres analytiques $\Delta = 73$ bits, $\alpha = 1{,}5$, $\beta = 1{,}5$,
$\lambda = 128$ (estimations param\'etriques sans mesure empirique contr\^ol\'ee),
l'heuristique de SOTP-High donne $0{,}45 \times 1{,}5 \times 7/128 \approx 0{,}037$.
Cette valeur n'est pas comparable au produit CLO des zk-SNARKs de conformit\'e
($\approx 0{,}572$), qui porte sur $\Relver$ et n'est pas une instance du
th\'eor\`eme (Remarque~\ref{rem:portee_cas_etude}).
\end{sloppypar}

% -----------------------------------------------------------------------
\subsection{Synth\`ese et s\'election de variante}
\label{subsec:synthese_selection}
% -----------------------------------------------------------------------
\begin{sloppypar}
\subsubsection{Comparaison des trois variantes}
\label{subsubsec:tableau_comparatif}

% -----------------------------------------------------------------------
\begin{table}[htbp]
\centering
\caption{Comparaison des variantes SOTP}
\label{tab:comparaison_sotp_construction}
\resizebox{\textwidth}{!}{%
\begin{tabular}{lp{4cm}p{4cm}p{4cm}}
\toprule
\textbf{Crit\`ere} & \textbf{SOTP-Low} & \textbf{SOTP-Base} & \textbf{SOTP-High} \\
\midrule
\multicolumn{4}{l}{\textit{\textbf{R\'egime op\'erationnel}}} \\
Entropie $H(\mathcal{C})$ & $< H^*(\lambda) \approx O(\sqrt{\lambda})$ & $H^*(\lambda) \leq H(\mathcal{C}) \leq \lambda/\log\lambda$ & $> \lambda/\log\lambda$ \\
Valeurs ($\lambda = 128$) & $< 11$ bits & $11$--$18$ bits & $> 18$ bits \\
\midrule
\multicolumn{4}{l}{\textit{\textbf{Techniques cryptographiques}}} \\
Couche Gold & PRG~\cite{BlumMicali84} & PRG~\cite{BlumMicali84} & ZK parfait (Groth16)~\cite{Groth16} \\
Couche Iron & Pedersen~\cite{Pedersen91} & Pedersen~\cite{Pedersen91} & knowledge soundness (q-SDH)~\cite{BonehBoyen08} \\
Couche Clay & $\phi:\mathcal{C}\to\{0,1\}^{H(\mathcal{C})}$ & JL $\Phi:\mathbb{R}^d\to\mathbb{R}^k$~\cite{JohnsonLindenstrauss84} & Poseidon + SNARK~\cite{Poseidon2019} \\
Primitives & SHA-256, Pedersen & Gaussienne, Pedersen & Groth16, Poseidon, BLS12-381~\cite{BLS12-381} \\
\midrule
\multicolumn{4}{l}{\textit{\textbf{M\'etriques CLO atteintes}}} \\
Confidentialit\'e & calculatoire (PRG) & calculatoire (PRG) & ZK parfait \\
$\Relver$ & $1 - \epsilon_{\text{DLog}}$ & $1 - \epsilon_{\text{DLog}}$ & $1 - \epsilon_{\text{q-SDH}}$ \\
$\Opp$ & d\'epend du contexte d'esp\`ece transmis & idem & idem \\
\midrule
\multicolumn{4}{l}{\textit{\textbf{Heuristique de conception (non prouv\'ee)}}} \\
Expression & $0{,}75 \cdot \Delta/\lambda$ & $0{,}5 \cdot \min\{\Delta/d, \Delta/(\beta\log_2 d)\}$ & $0{,}45 \cdot \beta\log_2\lambda/\lambda$ \\
Trusted setup & Non & Non & Oui~\cite{Bowe2017} \\
\midrule
\multicolumn{4}{l}{\textit{\textbf{Illustration (non comparable aux produits CLO)}}} \\
Syst\`eme de r\'ef\'erence & ECDSA~\cite{Johnson01} & Helios~\cite{Adida08} & zk-SNARKs conformit\'e~\cite{Groth16} \\
Valeur de l'heuristique & $\approx 0{,}205$ & $\approx 0{,}050$ & $\approx 0{,}037$ \\
Produit CLO (\'echelle $\Relver$) & $\approx 0{,}273$ & $\approx 0{,}437$ & $\approx 0{,}572$ \\
\bottomrule
\end{tabular}%
}
\end{table}
% -----------------------------------------------------------------------

L'association de chaque syst\`eme de r\'ef\'erence \`a une variante est
illustrative. Elle suit l'ordre relatif des entropies contextuelles de ces
syst\`emes (ECDSA, puis Helios, puis la conformit\'e financi\`ere). Les seuils
$H^*(\lambda)$ et $H^{**}(\lambda)$ de
l'Algorithme~\ref{alg:selection_variante_sotp} restent des crit\`eres de
dimensionnement \`a $\lambda$ fix\'e et ne servent pas \`a classer ces syst\`emes.

Les trois variantes forment un continuum adaptatif couvrant l'ensemble du spectre
d'entropie contextuelle. SOTP-Low privil\'egie la lisibilit\'e du contexte au prix d'une restriction sur
$H(\mathcal{C})$, id\'eal pour applications juridiques ou contractuelles~\cite{Mason07,Kerr04}. SOTP-Base
\'equilibre les trois m\'etriques via compression adaptative, polyvalent pour applications
gouvernementales et \'electorales~\cite{Adida08,Goodman15}. SOTP-High sacrifie l'opposabilit\'e pour la confidentialit\'e
parfaite et la scalabilit\'e, n\'ecessaire pour applications financi\`eres ou m\'edicales \`a
contexte tr\`es riche~\cite{Garratt15,BaselIII}.

\subsubsection{Crit\`ere de s\'election de variante}
\label{subsubsec:regle_selection}

\begin{algorithm}[H]
\caption{S\'election automatique de variante SOTP}
\label{alg:selection_variante_sotp}
\begin{algorithmic}[1]
\State \textbf{Entr\'ees} : $H(\mathcal{C})$, $\Delta$, $\lambda$, $\text{Priority}$,
  $\text{TrustedSetup\_OK}$
\State $H^*(\lambda) \leftarrow \sqrt{\lambda}$ ;\; $H^{**}(\lambda) \leftarrow \lambda/\log_2\lambda$
\If{$\text{TrustedSetup\_OK} = \texttt{False}$ \textbf{and} $H(\mathcal{C}) > H^{**}(\lambda)$}
  \State \textbf{ALERTE} : SOTP-High requis mais trusted setup indisponible~\cite{Bowe2017}
  \State \textbf{Recommander} : SOTP-Base avec perte d'opposabilit\'e, ou d\'eployer MPC
  \State \textbf{Retourner} SOTP-Base
\EndIf
\If{$H(\mathcal{C}) < H^*(\lambda)$}
  \State \textbf{Retourner} SOTP-Low
\ElsIf{$H(\mathcal{C}) \leq H^{**}(\lambda)$}
  \State \textbf{Retourner} SOTP-Base ;\; $d \leftarrow \lceil 4\Delta H(\mathcal{C})/\varepsilon^2 \rceil$
\Else
  \State \textbf{Retourner} SOTP-High ;\; compiler circuit $\mathcal{C}_{\text{circuit}}$
\EndIf
\end{algorithmic}
\end{algorithm}

Cet algorithme est un m\'eta-algorithme de d\'eploiement, distinct des
constructions cryptographiques. Il prend en entr\'ee les param\`etres
du contexte institutionnel et retourne la variante SOTP adapt\'ee, en
tenant compte des propri\'et\'es math\'ematiques du contexte
et des contraintes op\'erationnelles de l'environnement de
d\'eploiement.

La ligne~2 calcule les deux seuils qui partitionnent l'espace des
entropies contextuelles. $H^*(\lambda) = \sqrt{\lambda}$ est la
fronti\`ere inf\'erieure s\'eparant le r\'egime \`a contexte rare du r\'egime
de transition tandis que $H^{**}(\lambda) = \lambda/\log_2\lambda$ est la
fronti\`ere sup\'erieure s\'eparant le r\'egime de transition du r\'egime
\`a contexte riche. Ces deux seuils sont des fonctions croissantes de
$\lambda$, car plus le param\`etre de s\'ecurit\'e est \'elev\'e, plus les
fronti\`eres se d\'eplacent vers des entropies contextuelles plus grandes,
ce qui signifie qu'une m\^eme richesse contextuelle peut relever de
r\'egimes diff\'erents selon le niveau de s\'ecurit\'e cible.

Les lignes~3--7 constituent la garde de s\'ecurit\'e op\'erationnelle, qui
est la partie la plus importante de l'algorithme du point de vue du
d\'eploiement. SOTP-High requiert une \'etape de configuration
de confiance (\emph{trusted setup}) pour g\'en\'erer le CRS Groth16~;
si cette \'etape n'a pas \'et\'e conduite~(parce que
l'organisation n'a pas les ressources pour organiser un protocole
multi-parties, ou parce que le contexte r\'eglementaire ne l'autorise
pas) et que le contexte est n\'eanmoins riche ($H(\mathcal{C}) >
H^{**}(\lambda)$), l'algorithme \'emet une alerte explicite et ne
d\'egrade pas silencieusement la s\'ecurit\'e. Il recommande alors
soit d'accepter SOTP-Base avec la perte d'opposabilit\'e
associ\'ee au r\'egime pour lequel il n'est pas optimal, soit de d\'eployer
un protocole MPC pour r\'ealiser le trusted setup de mani\`ere
distribu\'ee. Ce traitement explicite de la d\'efaillance op\'erationnelle
distingue SOTP-High des deux autres variantes, qui n'ont aucune
d\'ependance envers une \'etape externe.

Les lignes~8--14 impl\'ementent la s\'election proprement dite par
comparaison en cascade. Si $H(\mathcal{C}) < H^*(\lambda)$, le contexte
est assez rare pour que l'encodage d\'eterministe de SOTP-Low soit
suffisant, sans qu'aucun param\`etre suppl\'ementaire ne soit n\'ecessaire. Si
$H(\mathcal{C})$ est compris entre les deux seuils, SOTP-Base est
s\'electionn\'e et la dimension de compression $d = \lceil 4\Delta
H(\mathcal{C})/\varepsilon^2 \rceil$ est calcul\'ee imm\'ediatement. Ce calcul est effectu\'e \`a la s\'election et non dans le \textsf{Setup}
pour permettre d'adapter $d$ \`a la valeur effective de $H(\mathcal{C})$
observ\'ee au moment du d\'eploiement. Si $H(\mathcal{C})$ exc\`ede
$H^{**}(\lambda)$ et que le trusted setup est disponible, SOTP-High
est s\'electionn\'e et le circuit arithm\'etique $\mathcal{C}_{\text{circuit}}$
est compil\'e. Cette compilation \'etant l'op\'eration la plus co\^uteuse
du processus, elle n'est d\'eclench\'ee qu'en
dernier recours.

La logique de s\'election refl\`ete ainsi une hi\'erarchie de pr\'ef\'erences
implicite. SOTP-Low est pr\'ef\'er\'e quand il est applicable (co\^ut
minimal, pas de d\'ependance externe), SOTP-Base est le choix par d\'efaut
en r\'egime interm\'ediaire (robuste, sans trusted setup), alors que SOTP-High
n'est d\'eploy\'e que lorsque la richesse contextuelle le rend
in\'evitable et que les conditions op\'erationnelles le permettent.

\subsubsection{M\'ethodologie de conception en cinq phases}
\label{subsubsec:methodologie_cinq_phases_sotp}

La s\'election algorithmique ci-dessus s'inscrit dans une m\'ethodologie plus large, organis\'ee en cinq phases.

\textbf{Phase 1. Analyse applicative.} Estimer $H(\mathcal{C})$ et identifier les
exigences minimales $\Conf_{\min}$, $\Rel_{\min}$, $\Opp_{\min}$.

\textbf{Phase 2. \'Evaluation institutionnelle.} Estimer $\Delta$ (capacit\'e de
d\'elib\'eration, expertise, m\'emoire institutionnelle) et $\alpha$ (granularit\'e des verdicts).
Ces estimations sont analytiques, car une mesure directe via protocoles cognitifs contr\^ol\'es
reste un travail futur.

\textbf{Phase 3. Test de faisabilit\'e.} V\'erifier l'in\'egalit\'e
\begin{equation}
\Conf_{\min} \cdot \mathsf{Rel}_{\mathrm{rec},\min} \cdot \Opp_{\min} \stackrel{?}{\leq}
\frac{1+\kappa/\lambda}{\lambda}\,\min\left\{1, \frac{\kappa/\lambda + \Delta}{\gamma}\right\}
\label{eq:test_faisabilite_sotp}
\end{equation}
pour les exigences exprim\'ees en fiabilit\'e de reconstruction
(\'equation~\eqref{eq:trilemme_cro_raffine}). Si l'in\'egalit\'e est viol\'ee, il est possible de
relaxer une exigence, d'augmenter $\Delta$ par investissement dans la
formation institutionnelle (effectif seulement si $\Delta < \gamma$) ou de r\'eviser la
granularit\'e des verdicts. Les exigences exprim\'ees en fiabilit\'e de v\'erification
$\Relver$ ne sont pas contraintes par ce test.

\textbf{Phase 4. S\'election et param\'etrage.} Appliquer l'Algorithme~\ref{alg:selection_variante_sotp}
et optimiser les primitives Gold/Iron/Clay dans le r\'egime identifi\'e.

\textbf{Phase 5. Preuves et validation.} \'Etablir les r\'eductions formelles aux
hypoth\`eses standard, puis mesurer $\Delta$ et l'opposabilit\'e obtenue selon les
protocoles du chapitre~\ref{chap:experimentation}.
\end{sloppypar}

% -----------------------------------------------------------------------
\subsection{Optimalit\'e universelle et analyse de sensibilit\'e}
\label{subsec:optimalite_sensibilite}
% -----------------------------------------------------------------------
\begin{sloppypar}
\subsubsection{Conjecture d'optimalit\'e universelle}
\label{subsubsec:optimalite_universelle}

\begin{conjecture}[Optimalit\'e universelle SOTP]\label{conj:optimalite_universelle}
\label{thm:optimalite_universelle}%
Sous les Hypoth\`eses~\ref{ass:extraction_bornee}--\ref{ass:regime}, pour chaque
r\'egime de richesse contextuelle, il existe une variante SOTP et une
constante $C_R \geq 1$ ind\'ependante de $\lambda$ telles que, pour tout
protocole CLO $\mathcal{P}'$ du m\^eme r\'egime,
\begin{equation}
\Conf(\mathcal{P}')\,\Relrec(\mathcal{P}')\,\Opp(\mathcal{P}')
\;\leq\; C_R \cdot \Conf(\text{SOTP})\,\Relrec(\text{SOTP})\,\Opp(\text{SOTP}).
\label{eq:optimalite_universelle}
\end{equation}
\end{conjecture}

\begin{remark}[Statut et argument heuristique]\label{rem:statut_optimalite}
La formulation initiale de cet \'enonc\'e, pr\'esent\'ee comme th\'eor\`eme sous la forme
$\text{produit}(\mathcal{P}') \leq c \cdot \text{produit}(\text{SOTP}) + o(1)$ avec
$c \in [0{,}45;\,0{,}75]$, a \'et\'e remplac\'ee par une formulation plus pr\'ecise, pour
quatre raisons. (i)~Avec $c < 1$ et $\mathcal{P}' = \text{SOTP}$, elle donnerait
$\text{produit} \leq 0{,}75 \times \text{produit}$, ce qui impose $C_R \geq 1$~;
(ii)~le terme additif $o(1)$ domine des produits tous en $O(1/\lambda)$~;
(iii)~son argument invoquait le trilemme sans l'Hypoth\`ese~\ref{ass:regime}~;
(iv)~elle s'appuyait sur les minorations des Remarques~\ref{rem:facteur_limitant_low},
\ref{rem:heuristique_sotp_base} et~\ref{rem:heuristique_sotp_high}, qui rel\`event
d'heuristiques de conception. La forme multiplicative ci-dessus, avec $C_R \geq 1$, est l'\'enonc\'e que
l'argument vise. Celui-ci proc\`ede en trois temps, qui consistent \`a majorer le produit de tout
$\mathcal{P}'$ par la forme raffin\'ee~\eqref{eq:trilemme_cro_raffine}~; minorer celui
de la variante SOTP du r\'egime~; former le ratio. La deuxi\`eme \'etape reste ouverte, car
elle requiert un mod\`ele dans lequel la couche Clay transmet un contexte d'esp\`ece
que l'institution ne d\'etient pas (perspective d\'evelopp\'ee au
chapitre~\ref{chap:resultats_discussions}). Les rapports indicatifs que donneraient
les heuristiques de conception, $1/0{,}75 \approx 1{,}33$ (SOTP-Low), $1/0{,}5 = 2$
(SOTP-Base) et $1/0{,}45 \approx 2{,}2$ (SOTP-High), fixent l'ordre de grandeur
de $C_R$ vis\'e.
\end{remark}

\begin{boiteArbitre}
L'arbitre comprend que la recherche du meilleur protocole peut
continuer. La borne lui dit ce qui est interdit~; savoir jusqu'o\`u les
constructions SOTP l'approchent reste une question ouverte, pos\'ee avec pr\'ecision.
\end{boiteArbitre}

\subsubsection{Analyse de sensibilit\'e}
\label{subsubsec:analyse_sensibilite}

L'analyse porte sur la borne raffin\'ee~\eqref{eq:trilemme_cro_raffine}, qui est le seul
\'enonc\'e prouv\'e, ainsi que sur les heuristiques de conception.

\textbf{Sensibilit\'e \`a $\Delta$.} Dans le r\'egime \`a contexte riche ($\Delta < \gamma$),
la borne est proportionnelle \`a $\kappa/\lambda + \Delta$. Augmenter $\Delta$ de $50$ \`a
$65$ bits (formation intensive des garants~\cite{Ericsson95}) rel\`eve la borne d'environ
$30\%$. Dans le r\'egime \`a contexte rare ($\Delta \geq \gamma$), la borne est celle de
Fano et ne d\'epend pas de $\Delta$.

\textbf{Sensibilit\'e \`a $\alpha$.} Dans le r\'egime \`a contexte riche, la borne varie en
$1/\gamma = 1/(\alpha H(\mathcal{C}) - 1)$, de sorte que des verdicts plus grossiers (passer de 12
niveaux \`a 6) rel\`event la borne et peuvent faire basculer l'institution dans le
r\'egime \`a contexte rare. Les institutions \`a verdicts tr\`es nuanc\'es (\'echelles de risque
fines, montants d'amendes gradu\'es) paient ainsi une p\'enalit\'e r\'eelle. L'optimum entre
pr\'ecision d\'ecisionnelle et compromis CRO d\'epend des priorit\'es applicatives.

\textbf{Sensibilit\'e \`a $\beta$.} La constante $\beta$ n'appara\^it pas dans la borne~;
elle intervient dans la garantie de confidentialit\'e du
Lemme~\ref{lem:confidentiality_guarantee}, o\`u un protocole plus r\'egulier (plus petit
$\beta$) resserre cette garantie. Dans l'heuristique de SOTP-High, qui est
proportionnelle \`a $\beta$, passer de $\beta = 1{,}5$ \`a $\beta = 1{,}2$ diminue
la valeur de $20\%$. $\beta$ est une propri\'et\'e des lois de preuve (divergence
KL) et ne d\'epend pas de la fonction de hachage utilis\'ee dans le circuit.
\end{sloppypar}

% -----------------------------------------------------------------------
\subsection{Extension quantique des constructions SOTP}
\label{subsec:extension_postquantique_sotp}
% -----------------------------------------------------------------------
\begin{sloppypar}
Les constructions SOTP pr\'esent\'ees dans cette section s'appliquent au r\'egime classique.
La section~\ref{sec:cro_pq} a \'etabli les bornes quantiques~; 
l'extension des constructions SOTP au r\'egime quantique requiert des adaptations 
propres \`a chaque variante.

\textbf{SOTP-Low post-quantique.} La couche Gold peut int\'egrer des PRG bas\'es sur 
les fonctions de hachage (dont l'algorithme de Grover~\cite{Grover96} r\'eduit la s\'ecurit\'e de $\lambda$
\`a $\lambda/2$ bits). L'engagement de Pedersen
reste parfaitement masquant, mais son caract\`ere liant repose sur DLog, que
l'algorithme de Shor r\'esout. Il doit donc \^etre remplac\'e par un engagement
post-quantique, par exemple \`a base de r\'eseaux sous SIS/LWE~\cite{Peikert2016}.

\textbf{SOTP-Base post-quantique.} La compression JL est ind\'ependante du mod\`ele 
adversaire~\cite{JohnsonLindenstrauss84,DasguptaGupta03}. Les adaptations concernent les couches Gold et Iron avec des primitives 
post-quantiques (engagement \`a base de r\'eseaux~\cite{Peikert2016}, signatures ML-DSA~\cite{Dilithium2022}).

\textbf{SOTP-High post-quantique.} Les zk-SNARKs Groth16 n\'ecessitent une attention 
particuli\`ere~\cite{Unruh2016,Chatterjee2021}. Les hypoth\`eses q-SDH et q-PKE peuvent \^etre affaiblies face \`a des 
adversaires quantiques (elles reposent sur des couplages, que l'algorithme de Shor
compromet). Les STARKs~\cite{BenSasson2019}, qui ne reposent que sur des fonctions de
hachage, constituent l'alternative post-quantique~; les Bulletproofs~\cite{Bulletproofs2018}
reposent sur DLog et ne sont pas post-quantiques.

Le coefficient de perte d'interpr\'etabilit\'e $\eta_q$
(D\'efinition~\ref{def:perte_interp_quantique}) s'applique \`a toutes les variantes
lorsque la preuve transite par un support quantique~:
\begin{equation}
\Opp^q \geq \Opp - \eta_q\,\frac{H(C)}{H(V|C)},
\end{equation}
et $\eta_q = 0$ lorsque la preuve reste une cha\^ine de bits classique.
\end{sloppypar}

% -----------------------------------------------------------------------
% Synth\`ese
% -----------------------------------------------------------------------
\begin{boiteArbitre}
L'arbitre referme le dossier SOTP. Il vient de terminer l'\'etape \textbf{S}. 
Cette section a modifi\'e sa mani\`ere de lire le Trilemme CRO. Une borne sup\'erieure
sans construction ressemble \`a une condamnation, alors qu'accompagn\'ee de trois constructions de
r\'ef\'erence elle devient une carte, dont une partie reste \`a lever.

Quatre contributions structurent ce qu'il retient. L'architecture commune
d'abord (Gold pour le masquage, Iron pour le liage, Clay pour lui), dont
la couche Clay s'adapte au r\'egime contextuel. Les trois constructions
ensuite, calibr\'ees chacune sur ce que son $\Delta$ lui permet d'extraire
sans lui en demander davantage, avec SOTP-Low quand le contexte est rare et
l'encodage d\'eterministe suffit, SOTP-Base quand la compression
Johnson-Lindenstrauss absorbe la transition, SOTP-High quand les preuves
succinctes Groth16 compriment un contexte arbitrairement riche en quelques
centaines d'octets. Viennent enfin la conjecture d'optimalit\'e universelle, qui
dit ce qu'il reste \`a d\'emontrer, puis la m\'ethodologie de s\'election, qui
transforme le r\'esultat en proc\'edure.

Il retient surtout la franchise du propos. Les propri\'et\'es cryptographiques des
constructions sont prouv\'ees~; les facteurs $0{,}75$, $0{,}5$ et $0{,}45$ sont des
heuristiques de conception que les syst\`emes qu'il conna\^it depuis le premier
chapitre ne servent pas \`a valider.

La certification formelle du trilemme, qui l'attend encore, fait l'objet de la section suivante.
\end{boiteArbitre}